% !TEX root = ../main.tex
\chapter{Determinants}
\label{ch:6}

We are halfway through the course, and we leave rectangular matrices behind. From here on the matrices are square, and the two big topics ahead are \emph{determinants} and \emph{eigenvalues}. The determinant comes first because the eigenvalue equation $\det(A-\lambda I)=0$ will lean on it directly.

The determinant is a single number $\det A$ --- also written $\abs{A}$ --- attached to every square matrix $A$. It is a remarkable amount of information squeezed into one number. The headline fact is the cleanest test for invertibility we will ever have: $A$ is invertible exactly when $\det A\ne 0$, and singular exactly when $\det A=0$. You already know one determinant from school,
\[
\det\begin{bmatrix} a & b \\ c & d\end{bmatrix}=ad-bc ,
\]
and the whole chapter is about understanding where that formula comes from, what it generalizes to, and what it is good for.

Strang's strategy --- and it is the right one --- is to resist writing down a big formula first. The determinant of an $n\times n$ matrix is a sum of $n!$ terms, and meeting that sum cold teaches you nothing. Instead we pin the determinant down by \emph{three properties} (think of them as axioms), and then \emph{derive} everything else: seven more properties, the big formula, the cofactor expansion, the inverse, Cramer's rule, and finally the geometric punchline that $\abs{\det A}$ is the volume of a box. Three rules in, a whole theory out.

\section{The Three Defining Properties}

Here are the three properties that \emph{define} the determinant. Every other fact in the chapter is a consequence of these.

\defn{The Determinant}{The \emph{determinant} is the unique function $\det$ that assigns a number to each square matrix and satisfies three properties:
\begin{itemize}[leftmargin=2.8em,labelwidth=2.2em,labelsep=0.6em,labelindent=0pt,align=left,itemsep=4pt]
   \item[\textbf{(P1)}] \textbf{The identity has determinant $1$:} $\det I = 1$.
   \item[\textbf{(P2)}] \textbf{Row exchange reverses the sign:} swapping any two rows of a matrix multiplies its determinant by $-1$.
   \item[\textbf{(P3)}] \textbf{The determinant is linear in each row separately.} Holding all other rows fixed, the determinant is a linear function of the one row that changes. This has two parts. Scaling one row by $t$ scales the determinant by $t$:
   \[
   \det\begin{bmatrix} t a & t b \\ c & d\end{bmatrix} = t\,\det\begin{bmatrix} a & b \\ c & d\end{bmatrix};
   \]
   and a sum in one row splits the determinant into a sum:
   \[
   \det\begin{bmatrix} a+a' & b+b' \\ c & d\end{bmatrix}
   = \det\begin{bmatrix} a & b \\ c & d\end{bmatrix}
   + \det\begin{bmatrix} a' & b' \\ c & d\end{bmatrix}.
   \]
\end{itemize}}

A word of warning about (P3), because it is the property people misread. Linearity is \emph{one row at a time}. It does \emph{not} say $\det(A+B)=\det A+\det B$ --- that is false (try two $2\times 2$ matrices). It says that if you change only the first row, leaving the second row alone, the determinant responds linearly to that one row. The other rows are frozen.

Two quick consequences of (P1) and (P2) are worth seeing right away. From (P1), $\det\!\left[\begin{smallmatrix}1&0\\0&1\end{smallmatrix}\right]=1$. Exchange the two rows and (P2) says
\[
\det\begin{bmatrix} 0 & 1 \\ 1 & 0\end{bmatrix} = -1 .
\]
So a permutation matrix --- an identity with its rows shuffled --- has determinant $+1$ or $-1$ depending on whether it took an \emph{even} or \emph{odd} number of row exchanges to build it. That sign is the seed of the whole theory.

\section{Seven Properties Derived from the Three}

Now the payoff. Properties \textbf{(P4)} through \textbf{(P10)} are not new assumptions --- each one follows from (P1), (P2), (P3). We prove them in order, because later ones use earlier ones.

\prop{(P4) Equal rows give determinant zero}{If two rows of $A$ are identical, then $\det A=0$.}
\pf{Suppose rows $i$ and $j$ are equal. Exchanging them changes nothing about the matrix, so the determinant is unchanged. But (P2) says a row exchange flips the sign, so the determinant equals its own negative: $\det A=-\det A$. Therefore $\det A=0$.}

\prop{(P5) Subtracting a multiple of one row from another doesn't change the determinant}{For $i\ne j$, subtracting $t$ times row $i$ from row $j$ leaves $\det A$ unchanged. These are exactly the row operations of elimination.}
\pf{Apply linearity (P3) to row $j$, which becomes $(\text{row }j) - t\,(\text{row }i)$:
\[
\det\begin{bmatrix} \vdots \\ \text{row } i \\ \vdots \\ \text{row } j - t\,\text{row } i \\ \vdots \end{bmatrix}
=\det\begin{bmatrix} \vdots \\ \text{row } i \\ \vdots \\ \text{row } j \\ \vdots \end{bmatrix}
- t\,\det\begin{bmatrix} \vdots \\ \text{row } i \\ \vdots \\ \text{row } i \\ \vdots \end{bmatrix}.
\]
The first term is $\det A$. The second matrix has row $i$ appearing twice (once in its own slot, once in slot $j$), so by (P4) its determinant is $0$. Hence the whole thing is $\det A$.}

This is the property that makes determinants \emph{computable}: the elimination steps that simplify a matrix do not change its determinant at all. In two dimensions the argument reads
\[
\det\begin{bmatrix} a & b \\ c-ta & d-tb\end{bmatrix}
=\det\begin{bmatrix} a & b \\ c & d\end{bmatrix}
- t\,\det\begin{bmatrix} a & b \\ a & b\end{bmatrix}
=\det\begin{bmatrix} a & b \\ c & d\end{bmatrix} - 0 .
\]

\prop{(P6) A zero row gives determinant zero}{If any row of $A$ is entirely zero, then $\det A=0$.}
\pf{The zero row is $0$ times itself, so by the scaling part of (P3) with $t=0$ we may pull out the factor $0$, giving $\det A = 0\cdot\det(\cdots)=0$.}

\prop{(P7) The determinant of a triangular matrix is the product of its diagonal entries}{If $A$ is upper triangular (or lower triangular) with diagonal entries $d_1,d_2,\dots,d_n$, then
\[
\det A = d_1 d_2 \cdots d_n .
\]}
\pf{First suppose every $d_i\ne 0$. Because $A$ is triangular, we can use the elimination operations of (P5) --- which do not change the determinant --- to clear out every off-diagonal entry, arriving at the diagonal matrix $\diag(d_1,\dots,d_n)$. Now factor $d_i$ out of row $i$ using the scaling part of (P3), one row at a time:
\[
\det\diag(d_1,\dots,d_n) = d_1 d_2 \cdots d_n \,\det I = d_1 d_2 \cdots d_n
\]
by (P1). If instead some $d_i=0$, then elimination produces a row of zeros before we finish, and (P6) gives $\det A = 0$, which is also the product $d_1\cdots d_n$ since one factor is zero. Either way the formula holds.}

This is the single most useful computational fact in the chapter. To find a determinant, run elimination down to a triangular form (the pivots sit on the diagonal), keep track of how many row exchanges you used, and multiply the pivots:
\[
\det A = (-1)^{(\#\text{ exchanges})}\,(\text{product of the pivots}) .
\]
That is exactly how a computer evaluates the determinant of a large matrix. It costs about $n^3/3$ operations --- nothing like the $n!$ terms of the big formula coming later.

\ex[The $2\times2$ formula falls out of elimination]{With $a\ne 0$, one elimination step clears the lower-left corner:
\[
\begin{bmatrix} a & b \\ c & d\end{bmatrix}
\;\longrightarrow\;
\begin{bmatrix} a & b \\ 0 & d-\tfrac{c}{a}b\end{bmatrix}.
\]
By (P5) the determinant is unchanged, and now the matrix is triangular, so by (P7)
\[
\det\begin{bmatrix} a & b \\ c & d\end{bmatrix}
= a\Bigl(d-\tfrac{c}{a}b\Bigr) = ad-bc . \quad\checkmark
\]
The pivots are $a$ and $d-\tfrac{c}{a}b$, and their product is the determinant.}

\prop{(P8) The determinant detects singularity}{$\det A=0$ if and only if $A$ is singular (not invertible). Equivalently, $\det A\ne 0$ exactly when $A$ is invertible.}
\pf{Run elimination on $A$. If $A$ is singular, elimination produces a row of zeros (a missing pivot), so by (P6) the determinant is $0$. If $A$ is nonsingular, elimination produces a full set of $n$ nonzero pivots $d_1,\dots,d_n$, and by (P7) the determinant is $\pm d_1 d_2\cdots d_n\ne 0$ (the sign records the row exchanges). So $\det A=0$ exactly in the singular case.}

This is the property that makes the determinant matter. It ties together every thread of the first half of the book: $\det A\ne 0$ is the same as ``the columns are independent,'' the same as ``$N(A)=\{0\}$,'' the same as ``$Ax=b$ has exactly one solution for every $b$,'' the same as ``$\rank A=n$'' (see Chapter~\ref{ch:4}).

\prop{(P9) The determinant of a product is the product of determinants}{For square matrices $A$ and $B$ of the same size,
\[
\det(AB) = (\det A)(\det B) .
\]}

We will not prove (P9) from scratch here --- the cleanest proof uses the $A=LU$ factorization, which we sketch for (P10) below --- but its consequences are immediate and constantly used.

\rmkb[Consequences of the product rule]{
A handful of corollaries come straight out of (P9):
\begin{itemize}
   \item \textbf{Inverse:} from $A^{-1}A=I$ and (P1), $(\det A^{-1})(\det A)=1$, so
   \[
   \det(A^{-1})=\frac{1}{\det A}\qquad(\text{when }A\text{ is invertible}).
   \]
   \item \textbf{Powers:} $\det(A^2)=(\det A)^2$, and in general $\det(A^k)=(\det A)^k$.
   \item \textbf{Scaling the whole matrix:} multiplying \emph{every} row by $t$ scales the determinant once per row, so for an $n\times n$ matrix
   \[
   \det(tA)=t^{\,n}\det A .
   \]
\end{itemize}
That last one is a volume statement in disguise: doubling all three edges of a box in $\RR^3$ multiplies its volume by $2^3=8$, which is exactly $\det(2A)=2^3\det A$.}

\prop{(P10) The determinant of the transpose equals the determinant}{For every square matrix, $\det(A\T)=\det A$.}
\pf{Use the factorization $A=LU$, where $L$ is lower triangular with $1$'s on its diagonal and $U$ is upper triangular with the pivots on its diagonal. By the product rule (P9), the claim $\det(A\T)=\det A$ becomes $\det(U\T L\T)=\det(LU)$, i.e.\ $\det(U\T)\det(L\T)=\det(L)\det(U)$. Now $L$ is triangular with diagonal $1$'s, so $\det L=\det(L\T)=1$ by (P7); and $U$ is triangular, so $\det(U\T)=\det U$ (transposing a triangular matrix keeps the same diagonal). Both sides are therefore equal to $\det U$. (When $A$ needs row exchanges, the same argument runs through $PA=LU$, and the permutation contributes $\pm 1$ on both sides.)}

The point of (P10) is leverage. The three defining properties were stated for \emph{rows}. Because transposing swaps rows and columns without changing the determinant, every row property automatically becomes a \emph{column} property. So: exchanging two columns flips the sign; the determinant is linear in each column; a repeated column or a zero column makes the determinant $0$; and so on. Rows and columns play identical roles.

\state{The ten properties at a glance}{
From the three axioms (P1)--(P3) we have derived everything:
\begin{itemize}[leftmargin=3.4em,labelwidth=2.8em,labelsep=0.6em,labelindent=0pt,align=left]
   \item[(P4)] equal rows $\Rightarrow$ determinant $0$;
   \item[(P5)] elimination steps don't change it;
   \item[(P6)] a zero row $\Rightarrow$ determinant $0$;
   \item[(P7)] triangular $\Rightarrow$ product of the diagonal;
   \item[(P8)] $\det A=0 \iff A$ singular;
   \item[(P9)] $\det(AB)=\det A\,\det B$;
   \item[(P10)] $\det(A\T)=\det A$ (so columns behave like rows).
\end{itemize}
Two of these are the workhorses: (P7) lets you \emph{compute} a determinant by elimination, and (P8) tells you what the answer \emph{means}.}

\rmkb[The one loose end]{There is a gap we have quietly stepped over. Property (P2) says a row exchange flips the sign, so the determinant of a permutation is $(-1)^{\#\text{exchanges}}$. But a given permutation can be reached by different numbers of exchanges --- could one route use seven swaps and another ten, giving both $-1$ and $+1$? If so the determinant would equal its own negative and the whole definition would collapse. It is a genuine theorem (the \emph{parity} of a permutation is well defined: an even permutation can never be reached by an odd number of swaps) that closes this gap. We take it as known; with it, the three properties really do define $\det$ uniquely.}

\section{The Big Formula: a Sum of \texorpdfstring{$n!$}{n!} Terms}

The three properties pin down the determinant, but so far our only recipe is ``run elimination.'' There is also an explicit closed formula, and deriving it is a beautiful application of linearity in the rows. We build it for $2\times 2$, watch the pattern in $3\times 3$, and then state it in general.

\subsection{Building the \texorpdfstring{$2\times 2$}{2x2} formula from scratch}

Take a general $2\times 2$ matrix and split each row into its two ``one-entry'' pieces using linearity (P3). The first row $(a,b)$ is $(a,0)+(0,b)$; the second row $(c,d)$ is $(c,0)+(0,d)$. Linearity in each row separately breaks the determinant into $2\times 2 = 4$ pieces:
\[
\det\begin{bmatrix} a & b \\ c & d\end{bmatrix}
=\det\begin{bmatrix} a & 0 \\ c & 0\end{bmatrix}
+\det\begin{bmatrix} a & 0 \\ 0 & d\end{bmatrix}
+\det\begin{bmatrix} 0 & b \\ c & 0\end{bmatrix}
+\det\begin{bmatrix} 0 & b \\ 0 & d\end{bmatrix}.
\]
Now look at the four pieces. The first and last each have a zero \emph{column}, so by (P6)-for-columns they vanish. Only the two pieces that pick \emph{one entry from each row and each column} survive:
\[
=\;\det\begin{bmatrix} a & 0 \\ 0 & d\end{bmatrix}
+\det\begin{bmatrix} 0 & b \\ c & 0\end{bmatrix}
\;=\; ad\,\det\!\begin{bmatrix} 1 & 0 \\ 0 & 1\end{bmatrix}
+ cb\,\det\!\begin{bmatrix} 0 & 1 \\ 1 & 0\end{bmatrix}
\;=\; ad - bc .
\]
The surviving determinants are determinants of permutation matrices: the diagonal one is $+1$, the off-diagonal one is $-1$. The number of survivors is exactly $2!=2$.

\subsection{The same idea in \texorpdfstring{$3\times 3$}{3x3}}

Splitting each of the three rows into its three one-entry pieces gives $3\times 3\times 3=27$ little determinants. A piece survives only if it has no repeated column --- i.e.\ it picks one entry from each row \emph{and} each column. The number of such choices is $3!=6$, and these are the six permutations of $(1,2,3)$. Each survivor is the chosen product times the sign ($\pm1$) of its permutation:
\[
\det\begin{bmatrix} a_{11} & a_{12} & a_{13} \\ a_{21} & a_{22} & a_{23} \\ a_{31} & a_{32} & a_{33}\end{bmatrix}
=\;
\begin{aligned}[t]
& a_{11}a_{22}a_{33} - a_{11}a_{23}a_{32} - a_{12}a_{21}a_{33} \\
&{}+ a_{12}a_{23}a_{31} + a_{13}a_{21}a_{32} - a_{13}a_{22}a_{31} .
\end{aligned}
\]
A handy mnemonic for $3\times 3$ (and \emph{only} $3\times 3$): the three products running ``down-and-to-the-right'' get a plus sign, the three running ``down-and-to-the-left'' get a minus sign. Do not try to extend this trick to $4\times 4$ --- it simply fails there.

\subsection{The general formula}

The pattern is now clear. For an $n\times n$ matrix there are $n!$ surviving terms, one for each permutation $(\alpha,\beta,\gamma,\dots,\omega)$ of $(1,2,\dots,n)$, each term taking one entry from every row and every column.

\thm{The Big Formula}{For an $n\times n$ matrix $A=(a_{ij})$,
\[
\det A = \sum_{\text{all } n! \text{ permutations}} \bigl(\pm\bigr)\, a_{1\alpha}\,a_{2\beta}\,a_{3\gamma}\cdots a_{n\omega},
\]
where $(\alpha,\beta,\gamma,\dots,\omega)$ runs over all permutations of the column indices $(1,2,\dots,n)$, and the sign is $+$ for an even permutation and $-$ for an odd one (the determinant of the corresponding permutation matrix).}

You can sanity-check the formula on the identity: every term contains a zero except the one from the trivial permutation $\alpha=1,\beta=2,\dots,\omega=n$, which contributes $a_{11}a_{22}\cdots a_{nn}=1$. So $\det I=1$, consistent with (P1). With more work one can verify (P2) and (P3) from the formula too; we take that on faith. The count $n!$ grows ferociously --- $2,6,24,120,\dots$ --- which is exactly why the elimination method (P7) is the one to use in practice.

\ex[Mostly zeros: spotting the surviving permutations]{Compute
\[
\det\begin{bmatrix} 0 & 0 & 1 & 1 \\ 0 & 1 & 1 & 0 \\ 1 & 1 & 0 & 0 \\ 1 & 0 & 0 & 1\end{bmatrix}.
\]
With so many zeros, almost every one of the $4!=24$ terms dies. A term survives only if it picks a nonzero entry from each row and column. Hunting them down:
\begin{itemize}
   \item The choice $a_{14}a_{23}a_{32}a_{41}=1\cdot1\cdot1\cdot1$ uses columns $(4,3,2,1)$. Turning $(4,3,2,1)$ back into $(1,2,3,4)$ takes \emph{two} exchanges (swap the outer pair, then the inner pair), so the sign is $+$: this term is $+1$.
   \item The choice $a_{13}a_{22}a_{31}a_{44}=1\cdot1\cdot1\cdot1$ uses columns $(3,2,1,4)$. That is \emph{one} exchange (swap positions $1$ and $3$) away from the identity, so the sign is $-$: this term is $-1$.
\end{itemize}
No other permutation avoids all the zeros. So $\det = (+1)+(-1)=0$. The matrix is singular, and indeed $(\text{row }1)-(\text{row }2)+(\text{row }3)-(\text{row }4)=0$ is a dependence among the rows.}

\section{Cofactors: Expanding Along a Row}

The big formula is correct but unwieldy. The \emph{cofactor expansion} reorganizes its $n!$ terms by factoring out the entries of one row, turning an $n\times n$ determinant into a combination of $(n-1)\times(n-1)$ determinants. It is the recursion that makes hand computation manageable, and it is the form we need for the inverse and for the eigenvalue polynomial later.

\subsection{The idea: group the big formula by the first row}

Go back to the $3\times 3$ big formula and gather all the terms containing $a_{11}$, then all containing $a_{12}$, then all containing $a_{13}$:
\[
\det A = a_{11}\underbrace{(a_{22}a_{33}-a_{23}a_{32})}_{C_{11}}
+ a_{12}\underbrace{(-a_{21}a_{33}+a_{23}a_{31})}_{C_{12}}
+ a_{13}\underbrace{(a_{21}a_{32}-a_{22}a_{31})}_{C_{13}} .
\]
Each parenthesis is (up to a sign) a $2\times 2$ determinant --- and not just any one. The factor multiplying $a_{11}$ uses only the entries \emph{not} in row $1$ or column $1$:
\[
C_{11}=+\det\begin{bmatrix} a_{22} & a_{23} \\ a_{32} & a_{33}\end{bmatrix},\quad
C_{12}=-\det\begin{bmatrix} a_{21} & a_{23} \\ a_{31} & a_{33}\end{bmatrix},\quad
C_{13}=+\det\begin{bmatrix} a_{21} & a_{22} \\ a_{31} & a_{32}\end{bmatrix}.
\]
The sign in front alternates $+,-,+$. This is the structure that holds in every dimension.

\defn{Minor and Cofactor}{Let $A$ be $n\times n$. The \emph{minor} $M_{ij}$ is the determinant of the $(n-1)\times(n-1)$ matrix obtained by deleting row $i$ and column $j$ of $A$. The \emph{cofactor} of $a_{ij}$ is the signed minor
\[
C_{ij}=(-1)^{i+j}\,M_{ij} .
\]
The sign $(-1)^{i+j}$ follows the checkerboard pattern
\[
\begin{bmatrix} + & - & + & \cdots \\ - & + & - & \cdots \\ + & - & + & \cdots \\ \vdots & & & \ddots\end{bmatrix},
\]
positive when $i+j$ is even, negative when $i+j$ is odd.}

\thm{Cofactor Expansion (Laplace Expansion)}{For any $n\times n$ matrix $A$, the determinant may be expanded along any row $i$:
\[
\det A = a_{i1}C_{i1}+a_{i2}C_{i2}+\cdots+a_{in}C_{in}=\sum_{j=1}^{n} a_{ij}C_{ij},
\]
or, because $\det A\T=\det A$, along any column $j$:
\[
\det A = a_{1j}C_{1j}+a_{2j}C_{2j}+\cdots+a_{nj}C_{nj}=\sum_{i=1}^{n} a_{ij}C_{ij}.
\]}
\pf{The cofactor $C_{ij}$ collects exactly the terms of the big formula that contain $a_{ij}$, with $a_{ij}$ divided out. Every term of the big formula contains exactly one entry from row $i$, so grouping all $n!$ terms by which row-$i$ entry they use partitions the sum into the $n$ groups $a_{ij}C_{ij}$. That the leftover factor in each group is $(-1)^{i+j}$ times the determinant of the matrix with row $i$ and column $j$ deleted is a bookkeeping check on the signs of the permutations; it works out to exactly the checkerboard sign. The column version follows by applying the row version to $A\T$.}

\rmkb[The two-by-two cofactor sanity check]{Expanding $\det\!\left[\begin{smallmatrix}a&b\\c&d\end{smallmatrix}\right]$ along the first row: the cofactor of $a$ is $+\det[d]=d$, and the cofactor of $b$ is $-\det[c]=-c$. So $\det = a\cdot d + b\cdot(-c)=ad-bc$, as it must be. The minus sign on the $b$ term is the $(-1)^{1+2}$ from the checkerboard.}

\subsection{Choose a good row or column}

The freedom to expand along \emph{any} row or column is the practical gift here: pick the one with the most zeros, so that most cofactors are never computed.

\ex[A tridiagonal recursion]{The $n\times n$ tridiagonal matrix of $1$'s --- with $1$'s on the diagonal and on the two adjacent diagonals, $0$'s elsewhere --- has a determinant that satisfies a tidy recursion. Write $D_n=\det A_n$. Expanding along the first row of
\[
A_4=\begin{bmatrix} 1 & 1 & 0 & 0 \\ 1 & 1 & 1 & 0 \\ 0 & 1 & 1 & 1 \\ 0 & 0 & 1 & 1\end{bmatrix}
\]
gives two surviving cofactors. The $a_{11}=1$ cofactor is $D_3$, the determinant of the lower-right $3\times 3$ tridiagonal block. The $a_{12}=1$ cofactor, after deleting row $1$ and column $2$ and applying the checkerboard minus sign, expands once more to leave $-1\cdot D_2$. So
\[
D_n = D_{n-1}-D_{n-2}.
\]
Starting from $D_1=1$ and $D_2=\det\!\left[\begin{smallmatrix}1&1\\1&1\end{smallmatrix}\right]=0$, the sequence is
\[
1,\;0,\;-1,\;-1,\;0,\;1,\;1,\;0,\;-1,\dots
\]
which repeats with period $6$. The determinant never grows --- a hint that these matrices hover near singularity.}

\section{The Inverse Formula}

Cofactors do more than compute a single number. Assembled correctly, they build the entire inverse $A^{-1}$. This is the explicit formula that the $2\times 2$ rule
\[
\begin{bmatrix} a & b \\ c & d\end{bmatrix}^{-1}
=\frac{1}{ad-bc}\begin{bmatrix} d & -b \\ -c & a\end{bmatrix}
\]
is secretly an instance of. Notice the ingredients on the right: the leading factor $1/\det A$, and a matrix of cofactors. The general pattern keeps both.

\defn{Cofactor Matrix}{The \emph{cofactor matrix} of $A$ is the $n\times n$ matrix $C$ whose $(i,j)$ entry is the cofactor $C_{ij}=(-1)^{i+j}M_{ij}$.}

\thm{The Cofactor Formula for the Inverse}{If $A$ is invertible, then
\[
A^{-1}=\frac{1}{\det A}\,C\T ,
\]
where $C$ is the cofactor matrix. \emph{Mind the transpose:} the cofactors of row $i$ of $A$ go into column $i$ of $A^{-1}$.}
\pf{We show $A\,C\T=(\det A)\,I$; dividing by $\det A$ then gives the formula. Look at the $(i,k)$ entry of $A\,C\T$. It is row $i$ of $A$ dotted with column $k$ of $C\T$ --- and column $k$ of $C\T$ is row $k$ of $C$, namely the cofactors $C_{k1},\dots,C_{kn}$:
\[
(A\,C\T)_{ik}=\sum_{j=1}^{n} a_{ij}\,C_{kj}.
\]
\emph{Diagonal entries ($i=k$).} The sum $\sum_j a_{ij}C_{ij}$ is exactly the cofactor expansion of $\det A$ along row $i$. So every diagonal entry equals $\det A$.

\emph{Off-diagonal entries ($i\ne k$).} The sum $\sum_j a_{ij}C_{kj}$ uses the \emph{entries} of row $i$ but the \emph{cofactors} of row $k$. That is precisely the cofactor expansion (along row $k$) of a modified matrix in which row $k$ has been replaced by a copy of row $i$. This modified matrix has two equal rows (rows $i$ and $k$ are both row $i$ of $A$), so by (P4) its determinant is $0$. Hence every off-diagonal entry is $0$.

Therefore $A\,C\T=(\det A)\,I$, and $A^{-1}=\frac{1}{\det A}C\T$.}

\ex[Back to the $2\times 2$ case]{For $A=\left[\begin{smallmatrix}a&b\\c&d\end{smallmatrix}\right]$ the cofactors are single entries with signs: $C_{11}=+d$, $C_{12}=-c$, $C_{21}=-b$, $C_{22}=+a$. So
\[
C=\begin{bmatrix} d & -c \\ -b & a\end{bmatrix},\qquad
C\T=\begin{bmatrix} d & -b \\ -c & a\end{bmatrix},\qquad
A^{-1}=\frac{1}{ad-bc}\begin{bmatrix} d & -b \\ -c & a\end{bmatrix}.
\]
The transpose is what swaps the off-diagonal cofactors into place. This is the familiar rule, now seen as the $n=2$ case of the cofactor formula.}

\rmkb[Why the cancellation works]{Strang's intuition: $\det A$ is a sum of products of $n$ entries, while each cofactor is a sum of products of $n-1$ entries. When $A$ multiplies $\tfrac{1}{\det A}C\T$, the missing entry on each side fills in so that the numerator becomes $\det A$ on the diagonal (giving $1$ after dividing) and a determinant-with-a-repeated-row off the diagonal (giving $0$). It is far easier to see this from the cofactor structure than from Gauss--Jordan elimination.}

\section{Cramer's Rule}

The inverse formula gives a startlingly explicit answer for the solution of $Ax=b$. When $A$ is invertible, $x=A^{-1}b=\frac{1}{\det A}C\T b$. Reading off the components of this product turns each unknown $x_j$ into a ratio of two determinants.

\thm{Cramer's Rule}{Let $A$ be invertible and $Ax=b$. Then for each $j$,
\[
x_j=\frac{\det B_j}{\det A},
\]
where $B_j$ is the matrix formed by replacing the $j$-th column of $A$ with the right-hand side $b$ (and leaving every other column of $A$ alone).}
\pf{From $x=\frac{1}{\det A}C\T b$, the $j$-th component is $x_j=\frac{1}{\det A}\sum_{i=1}^{n} C_{ij}\,b_i$ (the $j$-th row of $C\T$ is the $j$-th column of $C$, namely $C_{1j},\dots,C_{nj}$). Now the sum $\sum_i b_i C_{ij}$ is exactly the cofactor expansion, along column $j$, of the matrix whose column $j$ is $b$ and whose other columns are those of $A$ --- because the cofactors $C_{ij}$ depend only on the \emph{other} columns of $A$, which are untouched. That matrix is $B_j$, so $\sum_i b_i C_{ij}=\det B_j$, and $x_j=\det B_j/\det A$.}

\ex[Cramer on a $2\times 2$ system]{Solve
\[
\begin{bmatrix} 2 & 1 \\ 1 & 3\end{bmatrix}\begin{bmatrix} x_1 \\ x_2\end{bmatrix}=\begin{bmatrix} 3 \\ 5\end{bmatrix}.
\]
Here $\det A = 2\cdot 3-1\cdot 1=5$. Replace column $1$ by $b$ to get $B_1$, and column $2$ by $b$ to get $B_2$:
\[
\det B_1=\det\begin{bmatrix} 3 & 1 \\ 5 & 3\end{bmatrix}=9-5=4,\qquad
\det B_2=\det\begin{bmatrix} 2 & 3 \\ 1 & 5\end{bmatrix}=10-3=7.
\]
So $x_1=\tfrac{4}{5}$ and $x_2=\tfrac{7}{5}$. A quick check: $2\cdot\tfrac45+\tfrac75=\tfrac{8+7}{5}=3$ and $\tfrac45+3\cdot\tfrac75=\tfrac{4+21}{5}=5$. \checkmark}

\rmkb[Beautiful, but not for computing]{Cramer's rule is a gem of theory --- it writes each unknown as a clean ratio of determinants, and that explicit form is invaluable when you want to see \emph{how} the solution depends on the data. But as a numerical method it is terrible: evaluating $n+1$ determinants of size $n$ is wildly more work than the single elimination ($A=LU$, then back-substitute) that solves $Ax=b$ directly. Use Cramer's rule to understand, not to compute.}

\section{Determinant as Volume}

We close with the geometric meaning that justifies why this one number is worth all the trouble. The determinant is not just an algebraic test for invertibility --- it is, in absolute value, a \emph{volume}.

\thm{The Determinant Is a Volume}{Let $A$ be an $n\times n$ matrix. Then $\abs{\det A}$ is the volume of the $n$-dimensional box (parallelepiped) whose edges are the rows of $A$ --- equivalently the columns, since $\det A\T=\det A$. In $\RR^2$ this is the area of the parallelogram spanned by the two rows; in $\RR^3$, the volume of the parallelepiped spanned by the three rows.}

The clean way to believe this is to check that ``volume of the box'' obeys the same three defining properties (P1), (P2), (P3) that define the determinant --- and if two functions of a matrix obey the same three defining rules, they are the same function (up to the sign that the absolute value erases).

\begin{itemize}
   \item \textbf{(P1) for volume.} When $A=I$ the box is the unit cube, volume $1$. So volume agrees with $\det I=1$.
   \item \textbf{(P2) for volume.} Swapping two edges does not change the box, so it does not change the volume. The determinant flips sign, but $\abs{\det A}$ does not --- matched once we take absolute value. (More: for an orthogonal matrix $Q$ the box is a unit cube turned in space, volume $1$, and indeed $\det Q=\pm 1$.)
   \item \textbf{(P3a) for volume.} Doubling one edge of the box doubles its volume; scaling that edge by $t$ scales the volume by $\abs{t}$. This matches the scaling part of (P3).
   \item \textbf{(P3b) for volume.} The additive part is the only subtle one. It says that if one edge is split as a sum $v+v'$, the box on $v+v'$ has volume equal to the sum of the volumes of the boxes on $v$ and on $v'$ (the other edges shared). For a parallelogram this is a cut-and-paste fact: sliding along the shared base preserves area, and the two pieces reassemble correctly.
\end{itemize}

Since volume satisfies the same three rules, $\abs{\det A}$ \emph{is} the volume. Property (P4) is visible too: if two edges of the box coincide, the box is flat and has zero volume --- matching $\det A=0$ for a repeated row.

\rmkb[A geometric reading of singularity]{Now (P8) has a picture. $\det A=0$ means the box has collapsed to zero volume: the edge vectors (rows, or columns) lie in a lower-dimensional flat through the origin --- in $\RR^3$, squashed into a plane or a line. That is exactly the singular case from Chapter~\ref{ch:1}: dependent columns, a nontrivial null space, no inverse. Invertibility is the box having genuine $n$-dimensional volume.}

\ex[Areas of a parallelogram and a triangle]{The parallelogram with edge vectors $(a,b)$ and $(c,d)$ has area
\[
\abs{\det\begin{bmatrix} a & b \\ c & d\end{bmatrix}}=\abs{ad-bc}.
\]
The triangle formed by those same two edges is half the parallelogram, area $\tfrac12\abs{ad-bc}$. And a triangle with arbitrary vertices $(x_1,y_1)$, $(x_2,y_2)$, $(x_3,y_3)$ has area
\[
\frac12\left|\det\begin{bmatrix} x_1 & y_1 & 1 \\ x_2 & y_2 & 1 \\ x_3 & y_3 & 1\end{bmatrix}\right|,
\]
the extra column of $1$'s accounting for the fact that the vertices need not sit at the origin. Coordinates of corners in, area out --- the determinant is doing geometry.}

\begin{center}
\begin{tikzpicture}[scale=0.95,>=Stealth]
  % parallelogram spanned by u=(3,0.6) and v=(1,2)
  \coordinate (O) at (0,0);
  \coordinate (U) at (3,0.6);
  \coordinate (V) at (1,2);
  \coordinate (UV) at (4,2.6);
  \fill[blue!12] (O)--(U)--(UV)--(V)--cycle;
  \draw[->,very thick,blue!70!black] (O)--(U)
       node[midway,below right,font=\scriptsize] {$(a,b)$};
  \draw[->,very thick,red!70!black] (O)--(V)
       node[midway,above left,font=\scriptsize] {$(c,d)$};
  \draw[blue!70!black] (V)--(UV);
  \draw[red!70!black] (U)--(UV);
  \node[font=\scriptsize] at (2.0,1.15) {area $=\abs{ad-bc}$};
\end{tikzpicture}
\end{center}
\begin{center}
\small\textit{The parallelogram spanned by the two rows of $A$; its area is $\abs{\det A}$.}
\end{center}

\section{What to Carry Forward}

We built the determinant from three rules --- $\det I=1$, sign flip on row exchange, linearity in each row --- and out of them came the whole machine. Two derived properties carry the weight. Property (P7), the triangular product rule, makes the determinant \emph{computable}: eliminate, count exchanges, multiply the pivots. Property (P8), $\det A=0\iff A$ singular, makes it \emph{meaningful}: one number decides invertibility, independence, and solvability all at once.

From the same three rules we derived the big formula (a signed sum of $n!$ products, one per permutation), the cofactor expansion (the recursion down to smaller determinants), and from cofactors both the inverse $A^{-1}=\frac{1}{\det A}C\T$ and Cramer's rule $x_j=\det B_j/\det A$. These explicit formulas are theoretical treasures and computational dead ends --- when you actually need to solve a system, eliminate. Finally, the determinant is a volume: $\abs{\det A}$ measures the box on the rows (or columns) of $A$, and singular matrices are exactly the ones whose box has collapsed to zero volume.

\rmkb[Looking ahead]{The reason we needed determinants now is the next chapter. Eigenvalues are the numbers $\lambda$ for which $A-\lambda I$ is singular --- which by (P8) means
\[
\det(A-\lambda I)=0 .
\]
Expanding that determinant (cofactors will do it) produces the \emph{characteristic polynomial}, whose roots are the eigenvalues. Everything in Chapter~\ref{ch:7} starts here.}
